\documentclass[11pt]{article}

\usepackage[margin=1in]{geometry}
\usepackage{amsmath, amssymb, amsthm}
\usepackage{graphicx}
\usepackage{hyperref}
\usepackage{listings}
\usepackage{xcolor}

\title{MIT 6.s978 Deep Generative Models - Fall 2024 \\ Problem Set 4 Report}
\author{Your Name}
\date{\today}

\lstset{
    language=Python,
    basicstyle=\small\ttfamily,
    keywordstyle=\color{blue},
    stringstyle=\color{red},
    commentstyle=\color{green!60!black},
    numbers=left,
    numberstyle=\tiny\color{gray},
    stepnumber=1,
    numbersep=8pt,
    backgroundcolor=\color{white},
    showspaces=false,
    showstringspaces=false,
    showtabs=false,
    frame=single,
    tabsize=4,
    captionpos=b,
    breaklines=true,
    breakatwhitespace=false,
    escapeinside={\%*}{*)}
}

\begin{document}

\maketitle

\section*{Problem 1: DDPM on MNIST}

\subsection*{(a) Implementation of \texttt{ddpm\_schedules()}}
In this section, we implemented the forward noise scheduling parameters for the DDPM. 
The variance $\beta_t$ linearly interpolates between $\beta_{min}$ and $\beta_{max}$ over $T$ steps. 
From $\beta_t$, we computed $\alpha_t = 1 - \beta_t$ and the cumulative product $\bar{\alpha}_t = \prod_{i=1}^t \alpha_i$.
These parameters are essential for evaluating the forward diffusion process efficiently.

\begin{lstlisting}[caption={Implementation of \texttt{ddpm\_schedules()}}]
def ddpm_schedules(beta1, beta2, T):
    assert beta1 < beta2 < 1.0, "beta1 and beta2 must be in (0, 1)"

    beta_t = torch.linspace(beta1, beta2, T)
    
    alpha_t = 1.0 - beta_t
    alphabar_t = torch.cumprod(alpha_t, dim=0)
    
    oneover_sqrta = 1.0 / torch.sqrt(alpha_t)
    sqrt_beta_t = torch.sqrt(beta_t)
    sqrtab = torch.sqrt(alphabar_t)
    sqrtmab = torch.sqrt(1.0 - alphabar_t)
    mab_over_sqrtmab_inv = (1.0 - alpha_t) / sqrtmab

    return {
        "alpha_t": alpha_t,
        "oneover_sqrta": oneover_sqrta,
        "sqrt_beta_t": sqrt_beta_t,
        "alphabar_t": alphabar_t,
        "sqrtab": sqrtab,
        "sqrtmab": sqrtmab,
        "mab_over_sqrtmab": mab_over_sqrtmab_inv,
    }
\end{lstlisting}


\subsection*{(b) Implementation of \texttt{forward()}}
The forward process takes the clean image $x_0$ and time step $t$ to produce the noisy image $x_t$ using the formulation:
$x_t = \sqrt{\bar{\alpha}_t}x_0 + \sqrt{1 - \bar{\alpha}_t}\epsilon$, where $\epsilon \sim \mathcal{N}(0, I)$.
During training, we utilize Classifier-Free Guidance (CFG) by randomly dropping the condition $c$ based on the probability \texttt{drop\_prob} to allow the model to learn unconditional generation.

\begin{lstlisting}[caption={Implementation of \texttt{forward()}}]
    def forward(self, x, c):
        _ts = torch.randint(1, self.n_T + 1, (x.shape[0],)).to(self.device)
        noise = torch.randn_like(x)
        
        sqrtab = self.sqrtab[_ts - 1][:, None, None, None]
        sqrtmab = self.sqrtmab[_ts - 1][:, None, None, None]
        x_t = sqrtab * x + sqrtmab * noise
        
        context_mask = torch.bernoulli(torch.zeros_like(c).float() + 1. - self.drop_prob).to(self.device)
        return self.loss_mse(noise, self.nn_model(x_t, c, _ts / self.n_T, context_mask))
\end{lstlisting}


\subsection*{(c) Implementation of \texttt{sample()}}
The reverse sampling process progressively removes noise starting from pure Gaussian noise $x_T$. Using the Classifier-Free Guidance approach, we compute predictions for both conditioned and unconditioned states and extrapolate them over the guidance scale $w$.

\begin{lstlisting}[caption={Implementation of \texttt{sample()}}]
        for i in range(self.n_T, 0, -1):
            t_is = torch.tensor([i / self.n_T]).to(device).repeat(n_sample * 2)
            x_i_teased = x_i.repeat(2, 1, 1, 1)

            with torch.no_grad():
                eps = self.nn_model(x_i_teased, c_i, t_is, context_mask)

            eps1 = eps[:n_sample]
            eps2 = eps[n_sample:]
            guided_eps = (1.0 + guide_w) * eps1 - guide_w * eps2

            x_i = self.oneover_sqrta[i - 1] * (x_i - guided_eps * self.mab_over_sqrtmab[i - 1])

            if i > 1:
                z = torch.randn(n_sample, *size).to(device)
                x_i = x_i + self.sqrt_beta_t[i - 1] * z
        return x_i
\end{lstlisting}

\subsection*{(d) DDPM Model Training Results}
Please see the attached output plots for MNIST digits generated by the trained 400-timestep model.

\vspace{0.5cm}
\centerline{\textit{[Insert generated training image here]}}
\vspace{0.5cm}

\newpage
\section*{Problem 2: Ablation Study of DDPM on MNIST}

\subsection*{Varying Training Timesteps}
We varied the number of timesteps $T$ from the default $400$ to $2000$ and $500$.
\begin{itemize}
    \item \textbf{T = 500}: Offers a slightly longer degradation schedule than the baseline, with similar but slightly enhanced details.
    \item \textbf{T = 2000}: Substantially increases the resolution of Markov chain transitions, improving sample quality and robustness, but the training latency significantly increases.
\end{itemize}

\subsection*{Resampling Schedule for Inference}
To accelerate the test-time sampling process without retraining the model, we implemented a custom inference loop that skips intermediate sampling steps while correcting the parameter scale by mapping the original timestep schedules to the new reduced scale using the relation $\alpha_{t \to s} = \bar{\alpha}_t / \bar{\alpha}_s$. 
The model demonstrates adequate performance even with significantly reduced inference steps (50 and 100), effectively verifying the fast generative capability given DDIM-like schedule interpolation.

\begin{lstlisting}[caption={Implementation of \texttt{sample\_with\_resampling()}}]
def sample_with_resampling(model, n_sample, size, device, inference_steps, guide_w=0.0):
    # ... setup random noise and configurations
    step_indices = torch.linspace(model.n_T, 1, inference_steps).long()
    alphabar_t = model.alphabar_t
    
    for idx in range(len(step_indices)):
        i = step_indices[idx].item()
        
        cur_alphabar = alphabar_t[i - 1]
        prev_alphabar = alphabar_t[step_indices[idx + 1] - 1] if idx < len(step_indices) - 1 else torch.tensor(1.0).to(device)
        
        new_alpha = cur_alphabar / prev_alphabar
        new_beta = 1.0 - new_alpha
        oneover_sqrta = 1.0 / torch.sqrt(new_alpha)
        mab_over_sqrtmab = (1.0 - new_alpha) / torch.sqrt(1.0 - cur_alphabar)
        
        #... run eps prediction with standard CFG
        x_i = oneover_sqrta * (x_i - guided_eps * mab_over_sqrtmab)

        if idx < len(step_indices) - 1:
            z = torch.randn(n_sample, *size).to(device)
            x_i = x_i + torch.sqrt(new_beta) * z
    return x_i
\end{lstlisting}

\newpage
\section*{Problem 3: Continuous-time Stochastic Process}

\subsection*{(a) Derivation of $p(x(t))$}
Under the Variance Preserving SDE formulation, the parameters for the prior distribution degrade exponentially over the integral form of the noise schedule $B(t) = \int_0^t \beta(s)ds$. 
For our initial distribution consisting of a mixture of Gaussians $\sum \pi_i \mathcal{N}(x(0); \mu_i, \sigma_i)$, the analytical probability density at time $t$ remains a Gaussian Mixture:
$$ p(x(t)) = \sum_{i} \pi_i \mathcal{N}(x(t); \mu_{i, t}, \sigma^2_{i, t}) $$
where:
$$ \mu_{i, t} = e^{-\frac{1}{2}B(t)}\mu_i \quad \text{and} \quad \sigma^2_{i, t} = e^{-B(t)}\sigma_i^2 + 1 - e^{-B(t)} $$

\subsection*{(b) Calculation of $\nabla_{x(t)} \log p_t(x)$}
Using the probability density formulated above, the exact Score of the data distribution can be defined straightforwardly by chain rule:
$$ \nabla_x \log p_t(x) = \frac{\nabla_x p_t(x)}{p_t(x)} $$
Where $\nabla_x p_t(x)$ is simply a summation of derivations for each Gaussian component:
$$ \nabla_x p_t(x) = \sum_i \pi_i \left( -\frac{x - \mu_{i, t}}{\sigma^2_{i, t}} \right) \mathcal{N}(x(t); \mu_{i, t}, \sigma^2_{i, t}) $$

\subsection*{(c) SDE Simulation and Visualization}
\begin{lstlisting}[caption={Analytical Score matching code}]
    # Inside grad_log_p_xt(x_t, t)
    for i in range(len(means)):
        mean_i_t = mu_scale * means[i]
        var_i_t = (mu_scale ** 2) * (stds[i] ** 2) + var_t
        pdf_i = gaussian_pdf(x_t, mean_i_t, np.sqrt(var_i_t))
        
        p_xt_val += weights[i] * pdf_i
        grad_p_xt += weights[i] * pdf_i * (-(x_t - mean_i_t) / var_i_t)
        
    return grad_p_xt / (p_xt_val + 1e-10)
\end{lstlisting}

\vspace{0.5cm}
\centerline{\textit{[Insert generated SDE forward/reverse diagram here]}}
\vspace{0.5cm}
As visualized in the generated plot, the Forward SDE flattens out the bimodal distribution into standard Normal $\mathcal{N}(0, 1)$, while Reverse SDE recovers the original bimodal constraints perfectly referencing the mathematically exact Score calculation.

\end{document}